\documentclass[acmsmall]{acmart}

\setcopyright{none}

%% NOTE that a single column version is required for
%% submission and peer review. This can be done by changing
%% the \documentclass[...]{acmart} in this template to
%% \documentclass[manuscript,screen]{acmart}
%%
%% To ensure 100% compatibility, please check the white list of
%% approved LaTeX packages to be used with the Master Article Template at
%% https://www.acm.org/publications/taps/whitelist-of-latex-packages
%% before creating your document. The white list page provides
%% information on how to submit additional LaTeX packages for
%% review and adoption.
%% Fonts used in the template cannot be substituted; margin
%% adjustments are not allowed.
%%
%% \BibTeX command to typeset BibTeX logo in the docs
\AtBeginDocument{%
  \providecommand\BibTeX{{%
    \normalfont B\kern-0.5em{\scshape i\kern-0.25em b}\kern-0.8em\TeX}}}

%% Rights management information.  This information is sent to you
%% when you complete the rights form.  These commands have SAMPLE
%% values in them; it is your responsibility as an author to replace
%% the commands and values with those provided to you when you
%% complete the rights form.
%%% The following is specific to  and the paper
%%% 'Trace-Guided Synthesis of Effectful Test Generators'
%%% by Zhe Zhou, Ankush Desai, Benjamin Delaware, and Suresh Jagannathan.
%%%
\setcopyright{cc}
\setcctype{by}
\acmDOI{10.1145/3808264}
\acmYear{2026}
\acmJournal{PACMPL}
\acmVolume{10}
\acmNumber{PLDI}
\acmArticle{186}
\acmMonth{6}
\acmSubmissionID{pldi26main-p68-p}
\received{2025-11-13}
\received[accepted]{2026-04-03}

%%
%% For managing citations, it is recommended to use bibliography
%% files in BibTeX format.
%%
%% You can then either use BibTeX with the ACM-Reference-Format style,
%% or BibLaTeX with the acmnumeric or acmauthoryear styles, that include
%% support for advanced citation of software artifacts from the
%% biblatex-software package, also separately available on CTAN.
%%
%% Look at the sample-*-biblatex.tex files for templates showcasing
%% the biblatex styles.
%%

%%
%% The majority of ACM publications use numbered citations and
%% references.  The command \citestyle{authoryear} switches to the
%% "author year" style.
%%
%% If you are preparing content for an event
%% sponsored by ACM SIGGRAPH, you must use the "author year" style of
%% citations and references.
%% Uncommenting
%% the next command will enable that style.
%\citestyle{acmauthoryear}

%%
%% end of the preamble, start of the body of the document source.

\usepackage[normalem]{ulem}
\usepackage{changebar}
\usepackage{hyperref, wrapfig, amsthm, pict2e, booktabs, subcaption,
  subfiles, rotating, pifont, indentfirst, amsmath, amsfonts,
  mathrsfs, adjustbox, mathtools,booktabs, microtype, ebproof,
  enumitem, multirow, multicol, setspace, xcolor, lineno, listings,
  stmaryrd, xspace, stackengine, makecell}

\usepackage{empheq}
\newcommand*\widefbox[1]{\fbox{\hspace{2em}#1\hspace{2em}}}
\usepackage[tableposition=top]{caption}
\usepackage{tikz-cd}
\usetikzlibrary{decorations.pathmorphing}
\usepackage[ruled,linesnumbered,vlined]{algorithm2e}
\newcommand\mycommfont[1]{\footnotesize\ttfamily\textcolor{DeepGreen}{#1}}
\SetCommentSty{mycommfont}
\usepackage{graphicx}
\urlstyle{same}
\usepackage{minted}
\usemintedstyle{borland}
\usepackage{tikz}
\newcounter{markeq}
\setcounter{markeq}{0}
\usepackage{spacingtricks}
\newcommand*\tstrut[1]{\vstrut{#1}}
\newcommand*\bstrut[1]{\vstrut[#1]{0pt}}
\usepackage[T1]{fontenc}
\newcommand*\mystrut[1]{\vrule width0pt height0pt depth#1\relax}
% \usepackage{bm}

\newcommand{\tyName}{\mbox{\textsf{u}\textsc{Hat}}\xspace}

% ZZ's Macros
\usepackage{spacingtricks, commands, refinementtydef}

\begin{document}

%%
%% The "title" command has an optional parameter,
%% allowing the author to define a "short title" to be used in page headers.
\title{Trace-Guided Synthesis of Effectful Test Generators}

\author{Zhe Zhou}
\orcid{0000-0003-3900-7501}
\affiliation{%
  \institution{Purdue University}
  \city{West Lafayette}
  \country{USA}
}
\email{zhou956@purdue.edu}

\author{Ankush Desai}
\orcid{0000-0001-9006-0100}
\affiliation{%
  \institution{Snowflake}
  \city{San Mateo}
  \country{USA}
}
\email{ankushdesai@gmail.com}

\author{Benjamin Delaware}
\orcid{0000-0002-1016-6261}
\affiliation{%
  \institution{Purdue University}
  \city{West Lafayette}
  \country{USA}
}
\email{bendy@purdue.edu}

\author{Suresh Jagannathan}
\orcid{0000-0001-6871-2424}
\affiliation{%
  \institution{Purdue University}
  \city{West Lafayette}
  \country{USA}
}
\email{suresh@cs.purdue.edu}





%%
%% By default, the full list of authors will be used in the page
%% headers. Often, this list is too long, and will overlap
%% other information printed in the page headers. This command allows
%% the author to define a more concise list
%% of authors' names for this purpose.


%%
%% The abstract is a short summary of the work to be presented in the
%% article.


\begin{abstract}
  Several recently proposed program logics have incorporated notions
  of \emph{underapproximation} into their design, enabling them to
  reason about reachability rather than safety. In this paper, we
  explore how similar ideas can be integrated into an expressive type
  and effect system.  We use the resulting underapproximate type
  specifications to guide the synthesis of test generators that probe
  the behavior of effectful black-box systems. A key novelty of our
  type language is its ability to capture underapproximate behaviors
  of \emph{effectful} operations using symbolic traces that expose
  latent data and control dependencies, constraints that must be
  preserved by the test sequences the generator outputs.  We
  implement this approach in a tool called \name{}, and evaluate it on
  a diverse range of applications by integrating \name{}'s synthesized
  generators into property-based testing frameworks like QCheck and
  model-checking tools like P.  In both settings, the generators
  synthesized by \name{} are significantly more effective than the
  default testing strategy, and are competitive with state-of-the-art,
  handwritten solutions.
\end{abstract}


%%
%% The code below is generated by the tool at http://dl.acm.org/ccs.cfm.
%% Please copy and paste the code instead of the example below.
%%
\begin{CCSXML}
  <ccs2012>
     <concept>
         <concept_id>10011007.10011074.10011099.10011102.10011103</concept_id>
         <concept_desc>Software and its engineering~Software testing and debugging</concept_desc>
         <concept_significance>500</concept_significance>
         </concept>
   </ccs2012>
\end{CCSXML}

\ccsdesc[500]{Software and its engineering~Software testing and debugging}

%%
%% Keywords. The author(s) should pick words that accurately describe
%% the work being presented. Separate the keywords with commas.
\keywords{property-based testing, synthesis, type systems, underapproximation}

%%
%% This command processes the author and affiliation and title
%% information and builds the first part of the formatted document.
\maketitle

% \subfile{oopsla25/0-intro}
% \subfile{oopsla25/9-benchmark-MonkeyDB}
% \subfile{oopsla25/10-benchmark-Quickstorm}
% \subfile{oopsla25/11-benchmark-STLCFuzzer}

% \subfile{sections/reviews}
% \subfile{sections/rebuttal}

\subfile{sections/0-intro}
\subfile{sections/1-overview}
\subfile{sections/2-language}
\subfile{sections/3-algo}
\subfile{sections/4-evaluation}
\subfile{sections/5-related}

\section*{Acknowledgements}

We thank the anonymous reviewers for their detailed comments and
suggestions, and Prasita Mukherjee (Purdue), Zhendong Ang (NUS), Tingxuan Xia (PKU) for their helpful feedback on the artifact. This material is based upon work supported by the National Science Foundation under CCF-SHF 2321680.

\section*{Data Availability}
\label{sec:data}

An artifact containing our implementation, benchmark suite, and results is publicly available on Zenodo~\cite{artifact}.

%%
%% The next two lines define the bibliography style to be used, and
%% the bibliography file.

% \bibliographystyle{ACM-Reference-Format}
\bibliographystyle{plainnat}
% \bibliographystyle{plain}
% \bibliographystyle{ACM-Reference-Format}
% \citestyle{acmauthoryear}
\bibliography{bibliography}

\ifdefined\showappendix
\newpage
\appendix
\subfile{tech/outlines}
\subfile{tech/0-semantics}
\subfile{tech/1-typing}
\subfile{tech/2-denotation}
\subfile{tech/3-algo}
\subfile{tech/4-proof-0}
% \subfile{tech/4-proof-2}
% \subfile{tech/4-proof-3}
\subfile{tech/5-explanation}
\subfile{tech/6-evaluation}
\else
\fi

\end{document}
\endinput
%%
%% End of file `sample-acmsmall.tex'.
